% =============================================================================
% Methods section draft for the BoneStrengthML / V&V 40 MLOps paper
% -----------------------------------------------------------------------------
% Extracted from the replication package (branch: main, HEAD 27d3196,
% analysed on 2026-09-12). This file is meant to be \input into the paper's
% main .tex file; it contains no preamble. It relies only on standard
% packages (booktabs, amsmath, siunitx is NOT required).
%
% Conventions used in this draft:
%   * [TODO: ...]   information the authors must fill in / confirm.
%   * [CHECK: ...]  a claim the analyst could not verify from the repository.
% Every factual statement below is traceable to a file in the repository;
% the file is named in a LaTeX comment where useful.
%
% Section 8 ("Analyst notes") is NOT meant for the paper: it lists the
% misalignments between the abstract and the repository, and the open
% questions for the authors.
% =============================================================================

\section{Methods}
\label{sec:methods}

% -----------------------------------------------------------------------------
\subsection{Research design}
\label{sec:methods:design}

We conducted an engineering research study~[TODO: cite the ACM SIGSOFT
Empirical Standard for Engineering Research, or Wieringa's design science
framework] in which we (i)~designed and implemented an MLOps pipeline that
embeds verification-side credibility activities of ASME V\&V~40 as
first-class, configuration-driven pipeline stages, and (ii)~evaluated the
resulting artifact on a single, in-depth case study: BoneStrengthML, a
machine-learning surrogate of a finite-element (FE) bone-strain model.
The evaluation is exploratory and descriptive: we assess whether, and for
which credibility factors, the pipeline is able to generate, version, and make
queryable the evidence required by the model's context of use (COU).
The replication package (code, configuration, sample data, and this
document) is publicly available~[TODO: cite the GitHub repository / archived
DOI].

% -----------------------------------------------------------------------------
\subsection{Case study: BoneStrengthML and its context of use}
\label{sec:methods:case}

\subsubsection{The surrogate modelling task}
% Sources: README.md, docs/Context_of_use.md, config/BoneStrengthML.yml

BoneStrength is a resource-demanding, patient-specific finite-element
model that predicts the strain induced in the proximal femur during a
simulated side fall, starting from CT-derived femur characteristics.
BoneStrengthML is intended as a regression surrogate of BoneStrength for
research-grade fracture-risk assessment. Its input space has 95
dimensions:

\begin{itemize}
  \item 93 principal-component scores (\texttt{PC\_1}--\texttt{PC\_93}) of a
    morpho-densitometric statistical atlas of the femur, summarising the
    geometry and the mineral-density distribution of the bone sampled over a
    Cartesian lattice; and
  \item two integer angles, \texttt{PosAnt}~$\in[-30, 30]^\circ$ and
    \texttt{MedLat}~$\in[-30, 0]^\circ$, describing the orientation of the
    impact force during the fall.
\end{itemize}

The surrogate predicts two scalar outputs, each treated as an independent
single-output regression problem: \texttt{maxStrain\_11}, the maximum first
principal strain (tension), and \texttt{maxStrain\_33}, the absolute value of
the minimum third principal strain (compression). Both are dimensionless
strains, reported here in microstrain ($\mu\varepsilon$).
The admissible range of every input and output is declared in the pipeline
configuration (Section~\ref{sec:methods:config}) and is used both to validate
the training data and to sample the input space during verification.

\subsubsection{Dataset}
% Sources: data/raw/sampleDataset.csv (inspected), config/BoneStrengthML.yml,
% docs/Machine Learning Canvas.md (deleted in commit cbfd4f8, but informative)

All training and test data are synthetic: each row is one BoneStrength FE
simulation for one femur (atlas coordinates) under one fall orientation.
The dataset shipped with the replication package
(\texttt{data/raw/sampleDataset.csv}, semicolon-separated CSV, 72\,MB) has the
characteristics summarised in Table~\ref{tab:dataset}.
[CHECK: the file is called \emph{sample}Dataset; the code's docstrings mention
a $10^6 \times 97$ dataset. Confirm whether the experiments were run on this
file or on a larger one.]

\begin{table}[htbp]
  \centering
  \caption{Characteristics of the dataset in the replication package.}
  \label{tab:dataset}
  \begin{tabular}{lr}
    \toprule
    Rows (FE simulations)                          & 45\,332 \\
    Columns (95 inputs + 2 outputs)                & 97 \\
    Distinct femurs (distinct \texttt{PC\_*} tuples) & 1\,044 \\
    Distinct rows                                  & 37\,472 \\
    Distinct \texttt{PosAnt} values                & 61 (all integers in $[-30,30]$) \\
    Distinct \texttt{MedLat} values                & 31 (all integers in $[-30,0]$) \\
    \texttt{maxStrain\_11} range                   & 685--7\,365 $\mu\varepsilon$ \\
    \texttt{maxStrain\_33} range                   & 1\,657--19\,288 $\mu\varepsilon$ \\
    \bottomrule
  \end{tabular}
\end{table}

Rows for the same femur differ only in the two impact angles; consequently
the dataset is \emph{clustered by femur}, a property that drives the data
partitioning strategy (Section~\ref{sec:methods:split}). Exact duplicate rows
are present and are explicitly declared as admissible in the configuration
(\texttt{duplicate\_rows: true}).

\subsubsection{COU-derived credibility requirements}
\label{sec:methods:cou}
% Source: docs/Context_of_use.md

The COU document fixes the acceptance thresholds against which the surrogate
is assessed. They derive from the clinical fracture thresholds
($7300\,\mu\varepsilon$ in tension and $10\,400\,\mu\varepsilon$ in
compression)~[TODO: cite] and from the accuracy of the FE model being
surrogated (7\,\% RMSE normalised to the maximum value, i.e.\ an error of
$511\,\mu\varepsilon$ on \texttt{maxStrain\_11} and
$729\,\mu\varepsilon$ on \texttt{maxStrain\_33}). The surrogate is
required to be one order of magnitude more accurate than the FE model it
replaces, and the numerical (training-related) error is in turn required to
be one order of magnitude below the surrogate's accuracy threshold.
Because no error larger than the threshold is acceptable, accuracy is
assessed with the infinity norm of the error (maximum absolute error) rather
than with an average. Uncertainty quantification was deemed not applicable,
since the surrogate's inputs are simulation parameters with no associated
measurement error. Table~\ref{tab:cou} lists the requirements.

\begin{table}[htbp]
  \centering
  \caption{Credibility requirements derived from the context of use
    (\texttt{docs/Context\_of\_use.md}) and their encoding in the pipeline
    configuration (\texttt{config/BoneStrengthML.yml}). MAE$_\infty$ = maximum
    absolute error; MRE = mean relative error; MRER = maximum relative error
    ratio.}
  \label{tab:cou}
  \small
  \begin{tabular}{p{2.6cm}p{4.3cm}p{4.6cm}p{2.2cm}}
    \toprule
    Requirement & COU statement & Configuration & Status in code \\
    \midrule
    Accuracy (validation gate) &
      MAE$_\infty \le 51\,\mu\varepsilon$ (\texttt{maxStrain\_11}) and
      $\le 73\,\mu\varepsilon$ (\texttt{maxStrain\_33}) &
      \texttt{gate\_thresholds.validation}: MaximumAbsoluteError $5.1\times10^{-5}$ and $7.3\times10^{-5}$ &
      Declared; not enforced \\
    Input-space coverage &
      $\ge 10^6$ samples of the 95-D input space &
      Latin Hypercube sampling; \texttt{lhs\_samples} defaults to 1\,000 &
      Sampler implemented; requirement not met by default \\
    Convergence &
      relative error w.r.t.\ fully-trained model $\le 5\,\%$ &
      \texttt{convergence}: MRE $\le 0.05$, subsets $n \in \{100, 250, 500, 1000\}$ &
      Implemented, enabled \\
    Perturbation stability (smoothness) &
      input relative error not amplified; test 40\,\% of admissible space; perturb each input by $10^{-3}$ of its range &
      \texttt{smoothness}: MRER $\le 1$, \texttt{sample\_fraction} 0.4, \texttt{perturbation\_scaled\_magnitude} 0.001 &
      Implemented, enabled \\
    Numerical error &
      std.\ dev.\ of outputs over $\ge 10$ repeated trainings $\le 5.1\,\mu\varepsilon$ / $7.3\,\mu\varepsilon$ &
      \texttt{numerical\_error}: MaxStandardDeviation $5.1\times10^{-6}$ / $7.3\times10^{-6}$, \texttt{sets: 10} &
      Commented out; not implemented \\
    Existence / coverage check &
      (not in COU text) &
      \texttt{existence}: \texttt{samples: 1000} &
      Commented out; not implemented \\
    \bottomrule
  \end{tabular}
\end{table}

% -----------------------------------------------------------------------------
\subsection{Pipeline architecture and tooling}
\label{sec:methods:arch}
% Sources: pyproject.toml, uv.lock, docs/workflow-architecture.md

The pipeline is a Python~3.11 package (\texttt{bonestrength\_ml}) exposing a
command-line interface (\texttt{bsml}). It is organised in four layers
(Figure~\ref{fig:arch}~[TODO: adapt the architecture diagram from
\texttt{docs/workflow-architecture.md}]):

\begin{enumerate}
  \item a \textbf{configuration layer}, in which a single YAML file describes
    the dataset schema, the model-building procedure, the verification tests
    and the acceptance thresholds, and is validated at start-up against
    typed Pydantic models;
  \item a \textbf{library layer} with pure, orchestration-agnostic functions
    for data validation, splitting, training, metrics, verification, and
    experiment logging;
  \item an \textbf{orchestration layer}, where Prefect tasks wrap the library
    functions and Prefect flows compose them, providing retries, task
    naming, tagging, fan-out/fan-in parallelism and an optional monitoring
    UI; and
  \item a \textbf{tracking layer}, where MLflow records every trained model as
    a run, attaches verification evidence to the run of the selected model,
    and promotes it to the MLflow Model Registry when all verification
    tests pass.
\end{enumerate}

Table~\ref{tab:tools} lists the tools with the exact versions pinned in the
lock file of the replication package.

\begin{table}[htbp]
  \centering
  \caption{Software stack of the pipeline (versions from \texttt{uv.lock}).}
  \label{tab:tools}
  \small
  \begin{tabular}{llp{6.8cm}}
    \toprule
    Concern & Tool (version) & Role in the pipeline \\
    \midrule
    Language / packaging & Python 3.11; uv & Reproducible environment from a committed lock file \\
    Orchestration & Prefect 3.6.10 & Flows and tasks, retries, parallel task submission, run monitoring UI \\
    Experiment tracking & MLflow 3.8.1 & Run tracking (params, metrics, tags, artifacts), model logging with signature, Model Registry; local SQLite backend by default \\
    Configuration validation & Pydantic 2.12.5, PyYAML 6.0.3 & Typed schema for the YAML configuration; fail-fast on malformed configuration \\
    Data validation & Pandera 0.27.0, pandas 2.3.3 & DataFrame schema generated from the configuration (types, ranges, nullability, strict columns, duplicates) \\
    Model building & scikit-learn 1.8.0, XGBoost 3.1.3, CatBoost 1.2.10 & Estimators, cross-validated hyperparameter search, group-aware splitting \\
    Numerics & NumPy 2.4.0, SciPy 1.17.0 & Latin Hypercube sampling, metrics \\
    CLI / console & Typer 0.20.1, Rich 14.2.0 & \texttt{bsml train}, \texttt{bsml list-models}, \texttt{bsml mlflow-ui} \\
    Data versioning & --- & \emph{No data-versioning tool is present in the package (see Section~\ref{sec:notes}); the dataset is committed to Git.} \\
    \bottomrule
  \end{tabular}
\end{table}

\subsubsection{Configuration-driven design}
\label{sec:methods:config}
% Sources: config/BoneStrengthML.yml, bonestrength_ml/config.py

All experiment parameters live in one versioned YAML file
(\texttt{config/BoneStrengthML.yml}, about 630 lines) with four top-level
sections: \texttt{dataset} (name, path, provenance and version metadata, CSV
format metadata, and one entry per input and output field with its type,
admissible range, unit and nullability), \texttt{model\_building} (train/test
split, preprocessing flags, global hyperparameter-optimisation settings, and
the list of candidate models with their search spaces),
\texttt{verification} (the ordered list of verification tests, each with its
per-output thresholds and parameters), and \texttt{gate\_thresholds}
(validation acceptance criteria per output). The same field ranges thus serve
three purposes: data validation, Latin Hypercube sampling of the input space,
and perturbation sizing in the smoothness test. Because the configuration is
under version control together with the code, every pipeline execution is
tied to a specific revision of both.

% -----------------------------------------------------------------------------
\subsection{Experimental workflow}
\label{sec:methods:workflow}
% Sources: bonestrength_ml/workflows/flows.py, tasks.py

The main flow (\texttt{train\_all\_outputs}, invoked by \texttt{bsml train})
implements a \emph{train $\rightarrow$ compare $\rightarrow$ verify
$\rightarrow$ register} lifecycle (Figure~\ref{fig:workflow}~[TODO: adapt the
Mermaid flowchart from \texttt{docs/workflow-architecture.md}]). Its stages
are:

\begin{enumerate}
  \item \textbf{Set-up.} Load and validate the configuration; set the MLflow
    tracking URI and create or select the experiment.
  \item \textbf{Data loading and validation.} Read the CSV according to the
    declared format metadata and validate it against the schema generated
    from the configuration (Section~\ref{sec:methods:data}).
  \item \textbf{Partitioning.} Project the frame into features and targets and
    build one train/test split shared by all outputs
    (Section~\ref{sec:methods:split}).
  \item \textbf{Model building (fan-out).} Submit one Prefect task per
    (model family $\times$ output) pair, i.e.\ $8 \times 2 = 16$ tasks, each
    performing a cross-validated hyperparameter search and logging its own
    MLflow run as soon as it completes (Section~\ref{sec:methods:training}).
  \item \textbf{Comparison (fan-in).} Once all tasks have completed, log a
    summary MLflow run containing a comparison table of all models and
    select, for each output, the model with the lowest test RMSE
    (the \emph{winner}).
  \item \textbf{Verification (fan-out).} Run the configured verification
    tests on each winner in parallel and attach the results to the winner's
    MLflow run (Section~\ref{sec:methods:verification}).
  \item \textbf{Conditional registration.} Register a winner in the MLflow
    Model Registry if and only if every verification test passed
    (Section~\ref{sec:methods:gate}).
\end{enumerate}

Two reduced flows exist for iteration: one restricted to a single output
(steps 1--7 for that output) and one restricted to a single model family and
output (steps 1--4 only, with no verification and no registration).
Verification can also be skipped with a CLI flag, in which case no
registration takes place.

\subsubsection{Data loading and validation}
\label{sec:methods:data}
% Source: bonestrength_ml/data_loading.py

A Pandera \texttt{DataFrameSchema} is generated at run time from the
\texttt{dataset} section of the configuration: every declared column receives
a dtype (\texttt{Float64} or \texttt{Int64}, with coercion), a range check on
its admissible interval and a nullability constraint; the schema is
\emph{strict}, so undeclared columns cause a failure, and rows are required
to be jointly unique unless duplicates are explicitly allowed. Validation is
lazy, so all violations are collected and reported at once, and the task is
retried once on failure.

\subsubsection{Data partitioning}
\label{sec:methods:split}
% Source: bonestrength_ml/training/splitter.py

Because every femur appears in many rows (one per impact orientation), a
row-wise random split would place the same femur in both partitions and
inflate test performance. The pipeline therefore uses a \emph{group-aware}
split: a group label is derived from the 93 atlas coordinates (rows with
identical \texttt{PC\_*} values are the same femur) and
\texttt{GroupShuffleSplit} assigns 70\,\% of the \emph{femurs} to training
and 30\,\% to test, with a fixed seed. The same indices are reused for both
outputs, so all 16 models are trained and evaluated on identical samples, and
the same split function is invoked again by the verification tests, which
therefore see exactly the training partition used for model building. A
plain random row split remains available as a configuration option.

\subsubsection{Model building and hyperparameter optimisation}
\label{sec:methods:training}
% Sources: bonestrength_ml/training/model_factory.py, trainer.py, metrics.py,
% config/BoneStrengthML.yml

Eight regression families are compared (Table~\ref{tab:models}). For each,
the configuration lists the hyperparameters: scalar values are fixed, list
values define the search space. Search is cross-validated with 5 folds; when
the grouped split is active, folds are built with \texttt{GroupKFold} on the
same femur labels, so no femur is shared between a training and a validation
fold. The scoring function is the negative RMSE. Two search strategies are
available and selected per model: exhaustive \texttt{GridSearchCV} (default)
and \texttt{RandomizedSearchCV}, the latter capped at 50 sampled
configurations and used for the two largest grids (random forest and
XGBoost). The best estimator is refitted on the full training partition and
evaluated on both partitions with MSE, RMSE, MAE, maximum absolute error and
$R^2$. Gaussian-process regression is wrapped in a pipeline with a
\texttt{StandardScaler}; the other families receive raw inputs, since the
global \texttt{normalization}/\texttt{standardization} flags are set to
\texttt{false}. A single random seed (default 42) is propagated to the split,
the search, and every estimator that accepts one.

\begin{table}[htbp]
  \centering
  \caption{Candidate model families and hyperparameter search spaces
    (\texttt{config/BoneStrengthML.yml}). Sizes are the number of grid points;
    randomised search evaluates at most 50 of them.}
  \label{tab:models}
  \small
  \begin{tabular}{llp{6.2cm}rl}
    \toprule
    Family & Library & Searched hyperparameters & Grid size & Search \\
    \midrule
    Random forest & scikit-learn & n\_estimators (15), max\_depth (10), min\_samples\_split (3), min\_samples\_leaf (4) & 1\,800 & Randomised (50) \\
    XGBoost & XGBoost & n\_estimators (16), max\_depth (8), learning\_rate (5), subsample (5), colsample\_bytree (5), gamma (6) & 96\,000 & Randomised (50) \\
    Support-vector regression & scikit-learn & C (4), epsilon (7), gamma (2), kernel (3), degree (2) & 336 & Grid \\
    Linear regression & scikit-learn & --- (fit\_intercept fixed) & 1 & none \\
    Multi-layer perceptron & scikit-learn & hidden\_layer\_sizes (3), activation (2), alpha (5), learning\_rate\_init (5) & 150 & Grid \\
    Partial least squares & scikit-learn & n\_components (20) & 20 & Grid \\
    CatBoost & CatBoost & iterations (8), depth (3), learning\_rate (3), l2\_leaf\_reg (5), bagging\_temperature (6) & 2\,160 & Grid \\
    Gaussian process (scaled) & scikit-learn & kernel (3: RBF short, RBF unit, RBF+DotProduct+White), alpha (3) & 9 & Grid \\
    \bottomrule
  \end{tabular}
\end{table}

\subsubsection{Model selection}
% Source: bonestrength_ml/workflows/flows.py

For each output, the winner is the model with the lowest RMSE on the held-out
test partition. The comparison table (RMSE, $R^2$, MAE, maximum absolute
error for every model and output) and the winner per output are logged in a
dedicated \texttt{experiment\_summary} MLflow run.

\subsubsection{Verification suite}
\label{sec:methods:verification}
% Sources: bonestrength_ml/verification/*.py

Verification is implemented as a pluggable, registry-based suite: each test
is a module that registers a check function under a name; the orchestrator
iterates over the tests listed in the configuration, in the declared order,
and builds for each a context holding the winning model, the full feature and
target frames, the configuration entry and the seed. Every check returns a
uniform result (pass/fail flag, metrics, parameters, details), and the
results are aggregated into a report whose \texttt{all\_passed} property
gates registration. Adding a test requires a new module and a YAML entry
only; flows, logging and CLI are unchanged.

\paragraph{Input-space sampling.}
Both implemented tests evaluate the surrogate on synthetic points drawn by
Latin Hypercube Sampling (LHS) over the 95-dimensional box defined by the
configured admissible ranges, rather than on observed data. The sampler is a
basic (non-optimised) LHS with a fixed seed; a maximin variant is implemented
but not used. The number of points defaults to 1\,000 and can be raised per
test via the \texttt{lhs\_samples} parameter.

\paragraph{Convergence test.}
The test checks that the winner's predictions are stable with respect to the
amount of training data. For each subset size $k \in \{100, 250, 500,
1000\}$, $k$ rows are drawn at random from the training partition, a fresh
model with the winner's type and optimal hyperparameters is fitted on them,
and its predictions on the LHS points are compared with those of the winner
through the mean relative error, $\mathrm{MRE} = \frac{1}{n}\sum_i
|(\hat y^{(k)}_i - \hat y^{\star}_i)/\hat y^{\star}_i|$. The test passes if
the MRE is at or below the threshold (5\,\%) for the largest subset and for
every larger subset after the last exceedance; the smallest such $k$ is
reported as the minimum converging training-set size.

\paragraph{Smoothness (perturbation-stability) test.}
A fraction (40\,\%) of the LHS points is sampled; at each point and for each
input dimension $j$, the input is perturbed by $\pm\delta_j$, with $\delta_j$
equal to $10^{-3}$ of the admissible range of dimension $j$, and the
normalised central-difference partial derivative
\[
  \mathrm{MRER}_{ij} =
  \left|\frac{\hat y(x_i + \delta_j e_j) - \hat y(x_i - \delta_j e_j)}{2\,\delta_j}
  \cdot \frac{x_{ij}}{\hat y(x_i)}\right|
\]
is computed, i.e.\ the ratio between the relative change of the output and
the relative change of the input. The test passes if the maximum ratio over
all points and dimensions is at most 1, meaning that the surrogate does not
amplify relative input errors and its response surface shows no spurious
oscillation at the perturbation scale.

\paragraph{Tests specified but not yet implemented.}
The configuration contains, commented out, a \emph{numerical error} test
(train the optimal configuration on 10 independent random subsets, require
the maximum per-point standard deviation of the predictions to stay below
$5.1\,\mu\varepsilon$ / $7.3\,\mu\varepsilon$) and an \emph{existence} test
(1\,000 samples, no threshold). Reference skeletons of both convergence and
numerical-error tests written by the FE-model developers are included in
\texttt{references/} and were the basis for the pipeline implementations.

\subsubsection{Registration gate}
\label{sec:methods:gate}
% Sources: bonestrength_ml/workflows/tasks.py, training/mlflow_utils.py

A winner is registered in the MLflow Model Registry under the name
\texttt{BoneStrengthML\_\{output\}\_\{model\}} only if every configured
verification test passed. A failed test leaves the model logged and
inspectable, with its failure recorded, but not promoted.

% -----------------------------------------------------------------------------
\subsection{Evidence logging and traceability}
\label{sec:methods:evidence}
% Source: bonestrength_ml/training/mlflow_utils.py

MLflow is used as the evidence store. The pipeline produces three kinds of
runs within one experiment:

\begin{itemize}
  \item \textbf{One run per trained model} (16 per full execution), created by
    the training task itself as soon as the model is fitted. It records the
    model type and label, the target output, the split settings, the
    cross-validation settings, the selected hyperparameters (prefixed
    \texttt{best\_}), train and test metrics (\texttt{train\_*},
    \texttt{test\_*}), the fitted estimator with an inferred input/output
    signature and an input example, and the complete cross-validation table
    as a CSV artifact.
  \item \textbf{One summary run} per execution with the model-comparison
    table and the best model and RMSE per output.
  \item \textbf{Verification evidence attached to the winner's run}: for each
    test, metrics (\texttt{verification.<test>.<metric>}), parameters
    (thresholds, subset sizes, sample counts), a pass/fail tag
    (\texttt{verification.<test>}), an overall
    \texttt{verification.status} tag, and the full structured report as a
    JSON artifact (\texttt{verification/report.json}).
\end{itemize}

Because thresholds, measured values and outcomes are logged as MLflow
parameters, metrics and tags, the evidence is queryable through the MLflow
search API and UI (e.g.\ all runs with \texttt{verification.status =
passed}), and the registered model version points back to the run holding
its evidence. Prefect adds an execution-level trace (task names such as
\texttt{train\_<model>\_<output>}, tags \texttt{training} and
\texttt{verification}, retries and timings) when a Prefect server is used.

% -----------------------------------------------------------------------------
\subsection{Reproducibility and execution environment}
\label{sec:methods:repro}

The environment is fully pinned by a committed \texttt{uv.lock}
(Table~\ref{tab:tools}); all stochastic components (train/test split,
randomised search, estimators, LHS, subset sampling) derive from a single
seed given on the command line (default 42). The training task is retried up
to twice on failure and the data-loading task once. Hyperparameter search
runs single-threaded by default (\texttt{n\_jobs: 1}); parallelism is
obtained at the Prefect level, where the 16 training tasks are submitted
concurrently to a thread-pool task runner. [TODO: report the hardware, the
wall-clock time of a full execution, and the Prefect/MLflow deployment used
for the reported results (local SQLite by default; a server deployment guide
for Ubuntu is in \texttt{docs/DEPLOYMENT.md}).]

% -----------------------------------------------------------------------------
\subsection{Mapping to ASME V\&V 40 verification-side credibility factors}
\label{sec:methods:vv40}

[CHECK: the mapping below is the analyst's reconstruction from the code and
the COU document; the repository itself only states that the suite is
``inspired by ASME V\&V 40''. Please confirm the factor names and the
credibility levels you intend to claim.]

\begin{table}[htbp]
  \centering
  \caption{Tentative mapping between V\&V 40 verification-side activities and
    pipeline stages.}
  \label{tab:vv40}
  \small
  \begin{tabular}{p{3.4cm}p{5.2cm}p{4.6cm}}
    \toprule
    V\&V 40 factor (analogue for an ML surrogate) & Pipeline evidence & Automation status \\
    \midrule
    Code verification -- software quality assurance &
      Typed configuration validation; strict data-schema validation; versioned code, configuration and lock file; pluggable test registry &
      Automated \\
    Code verification -- numerical code verification &
      Reference skeletons from FE-model developers vs.\ pipeline implementation (\texttt{references/}) &
      Manual (no automated tests in the package) \\
    Calculation verification -- discretisation error (analogue: training-set size) &
      Convergence test: MRE vs.\ subset size, minimum converging $k$ &
      Automated, enabled \\
    Calculation verification -- numerical solver error (analogue: training stochasticity) &
      Numerical-error test: per-point std.\ dev.\ over repeated trainings &
      Specified, not implemented \\
    Calculation verification -- use error / input sensitivity (analogue: perturbation stability) &
      Smoothness test: normalised partial derivatives on perturbed LHS points &
      Automated, enabled \\
    Input-space coverage &
      LHS over configured admissible ranges &
      Automated; sample size below COU requirement \\
    Accuracy (validation-side, listed for completeness) &
      \texttt{gate\_thresholds.validation} (MAE$_\infty$) &
      Declared, not enforced \\
    \bottomrule
  \end{tabular}
\end{table}

% =============================================================================
\section*{Analyst notes: misalignments with the abstract and open questions}
\label{sec:notes}
% Not part of the paper. Ordered by relevance to the abstract's claims.

\paragraph{N1 -- DVC is not used.}
The abstract states that the pipeline ``integrates DVC for data versioning''.
The repository contains no \texttt{.dvc/} directory, no \texttt{*.dvc}
pointer files, no \texttt{dvc.yaml}/\texttt{dvc.lock}, and DVC is not a
dependency in \texttt{pyproject.toml} or \texttt{uv.lock}. The 72\,MB dataset
is committed directly to Git. Data provenance is captured only by free-text
fields in the configuration (\texttt{version: v1};
\texttt{provenance: "<DOI or internal reference>"}, still a placeholder), and
these fields are not logged to MLflow. Either DVC must be added (and the
data-version hash logged with each run), or the abstract should be reworded
(e.g.\ ``Git-versioned configuration and data'').

\paragraph{N2 -- The accuracy threshold is declared but never enforced.}
The abstract lists ``accuracy'' among the COU-derived thresholds the surrogate
is assessed against. The configuration declares
\texttt{gate\_thresholds.validation} (maximum absolute error
$5.1\times10^{-5}$ and $7.3\times10^{-5}$), and the maximum absolute error is
computed and logged as \texttt{test\_max\_ae}, but no code reads
\texttt{gate\_thresholds}: winners are chosen by lowest test RMSE and
registration is gated only on the verification tests. A model violating the
COU accuracy requirement can therefore be registered. Question: is the
validation gate intentionally left as a human-in-the-loop step (which would
fit the abstract's conclusion), or is it a missing feature?

\paragraph{N3 -- Numerical-error test is not implemented.}
The abstract names ``numerical error'' as one of the assessed thresholds; the
COU requires it. The test is commented out in the configuration, has no
implementation in \texttt{bonestrength\_ml/verification/}, and its metric
name (\texttt{MaxStandardDeviation}) exists only in the Pydantic enum. A
reference skeleton exists in \texttt{references/TestNumericalError.py}, and
the stale branch \texttt{origin/vv4ml-test} (June 2026, author CittiX) calls
\texttt{vv4ml.run\_numerical\_error\_test} from an external package
\texttt{vv4ml} that is not in this repository. Question: is \texttt{vv4ml} a
separate library under development that will replace the in-package
verification suite? If so, the paper should say which one is evaluated.

\paragraph{N4 -- Input-space coverage falls short of the COU.}
The COU requires at least $10^6$ samples of the input space; both tests
default to 1\,000 LHS points (400 of which are used by the smoothness test),
and the committed configuration does not override \texttt{lhs\_samples}. The
commented-out \texttt{existence} test also uses 1\,000. Either the runs
reported in the paper set \texttt{lhs\_samples} much higher (please confirm
and document the value), or the ``input-space coverage'' factor should be
reported as not met at the required level.

\paragraph{N5 -- Convergence subsets are small relative to the training set.}
With about 70\,\% of 45\,332 rows in training ($\approx$ 31\,000 rows), the
largest convergence subset (1\,000 rows) is about 3\,\% of the training
partition. ``Converged'' therefore means that a model trained on 1\,000 rows
is within 5\,\% MRE of the winner, not that the error curve has flattened
near the full size. Consider adding larger subsets (e.g.\ 5\,000, 10\,000,
20\,000) or stating this limitation.

\paragraph{N6 -- Traceability gaps in the logged evidence.}
For a claim of a ``complete, traceable evidence chain'', the following are
not logged to MLflow at present: the random seed, the dataset name/version/
path (or a content hash), the configuration file (or its hash), and the
Prefect flow-run id. MLflow does record the Git commit automatically when
run from a repository, which partially covers code and configuration.
Additionally, the per-run parameters \texttt{cv\_folds} and
\texttt{optimization\_method} are taken from the \emph{global} optimisation
block, so runs for random forest and XGBoost are logged as
\texttt{GridSearch} although they use \texttt{RandomizedSearchCV}
(\texttt{mlflow\_utils.py}, \texttt{log\_training\_result}).

\paragraph{N7 -- Model selection on the test partition.}
The winner is the model with the lowest RMSE on the same held-out partition
whose metrics are then reported; there is no separate validation partition
(cross-validation is used only inside each family's search). This is a
common but debatable choice for a surrogate that will be assessed against
absolute thresholds; a reviewer may ask for selection on cross-validation
score instead.

\paragraph{N8 -- Configuration options that are parsed but inert.}
\texttt{preprocessing.normalization}, \texttt{preprocessing.standardization},
\texttt{optimization.feature\_selection} and \texttt{optimization.scoring}
are validated by Pydantic but have no effect (scoring is always RMSE; only
the GPR pipeline standardises inputs). The configuration also still contains
Italian design notes as comments. These should be removed or implemented
before the package is archived.

\paragraph{N9 -- No automated tests, no CI.}
The package has no \texttt{tests/} directory and no CI configuration; the
only executable check of the verification code is the comparison with the
reference skeletons. This matters for the ``software quality assurance''
row of Table~\ref{tab:vv40}.

\paragraph{N10 -- Results are not in the package.}
\texttt{mlruns/} and \texttt{mlruns.db} are git-ignored, so the evidence
chain produced for the paper is not in the replication package. An exported
MLflow experiment (or at least the \texttt{report.json} artifacts and the
comparison table) should be archived alongside the code.

\paragraph{N11 -- Dataset identity.}
The file is named \texttt{sampleDataset.csv}, the code docstrings mention a
$10^6$-row dataset, and the notebook comments say validation is ``expected to
fail if config ranges don't match data''. Please confirm which dataset the
reported experiments used, whether it is the published one referred to in
the configuration's \texttt{source} field, and its citation.

\paragraph{N12 -- Unmerged parallelism branch.}
Branch \texttt{origin/optimize-training-cpu-parallelism} (June 2026) adds a
bounded CPU budget for the hyperparameter search. If the reported runs used
it, the execution-environment paragraph should say so and the branch should
be merged before archiving.

\paragraph{Open questions for the authors (summary).}
\begin{enumerate}
  \item Which V\&V 40 factors and credibility levels do you intend to claim
    (to fill ``[N of M]'' in the abstract), and is Table~\ref{tab:vv40} the
    mapping you have in mind?
  \item Is the accuracy gate intentionally manual?
  \item Will numerical-error and existence tests be implemented in this
    package or in \texttt{vv4ml}?
  \item What \texttt{lhs\_samples} value was used for the reported results?
  \item Was DVC used in an environment outside this repository, or should it
    be dropped from the abstract?
  \item Which dataset and hardware were used, and how long did a full
    execution take?
\end{enumerate}
